(a + b)^n = \sum_{r=0}^{n} \binom{n}{r} a^{n-r} b^r

General term:  T_{r+1} = \binom{n}{r} a^{n-r} b^r

Number of terms in  (a+b)^n:  n+1

Sum of all coefficients: set  a = b = 1 \rightarrow 2^n

Sum of numerical coefficients in  \left(x + \frac{1}{x}\right)^5: set  x = 1 \rightarrow 2^5 = 32

Specific term ---  (a+b)^{22}, term  ka^wb^{20}:
r = 20 \rightarrow T_{21} = \binom{22}{20} a^2 b^{20} = 231\,a^2 b^{20}
k = 231,\; w = 2 \rightarrow k + w = 233

Equal consecutive terms condition:
T_{r+1} = T_{r+2} \rightarrow \binom{n}{r} = \binom{n}{r+1}
\rightarrow r+1 = n-r \rightarrow r = \frac{n-1}{2}

Example (2023 \#11): 5th and 6th terms of  (x+y)^n  equal
\binom{n}{4} = \binom{n}{5} \rightarrow 5 = n - 4 \rightarrow n = 9
3rd term:  \binom{9}{2}x^7y^2 = 36x^7y^2 \rightarrow coefficient = 36

State 2022 \#9:  (2x+c)^8, 4th and 5th terms equal (numerical coefficients)
T_4 = \binom{8}{3}(2x)^5 c^3 = 56 \cdot 32\, x^5 c^3
T_5 = \binom{8}{4}(2x)^4 c^4 = 70 \cdot 16\, x^4 c^4
56 \cdot 32 = 70 \cdot 16 \cdot c
c = \frac{56 \cdot 32}{70 \cdot 16} = \frac{1792}{1120} = \frac{8}{5}

Large exponent (2020 \#17):  (2x+7)^{101}, coefficient of  x^{53}
T_{54} = \binom{101}{53} \cdot 2^{53} \cdot x^{53} \cdot 7^{48}

Pascal's Triangle shortcut:  \binom{n}{r} + \binom{n}{r+1} = \binom{n+1}{r+1}
